% ============================================================================
%  Election-night model: specification
%  Precinct x vote-mode model for live returns.
%  Compile: pdflatex tx_model_spec.tex
% ============================================================================
\documentclass[11pt]{article}

\usepackage[margin=1in]{geometry}
\usepackage{amsmath, amssymb}
\usepackage{microtype}
\usepackage{enumitem}
\usepackage[colorlinks=true, allcolors=blue]{hyperref}

\setlist[itemize]{itemsep=2pt, topsep=4pt}
\newcommand{\dem}{\mathrm{dem}}
\newcommand{\rep}{\mathrm{rep}}
\newcommand{\oth}{\mathrm{oth}}
\newcommand{\refm}{\mathrm{ref}}
\newcommand{\logit}{\operatorname{logit}}
\newcommand{\softmax}{\operatorname{softmax}}
\newcommand{\ICC}{\mathrm{ICC}}
\newcommand{\var}[1]{\texttt{#1}}

\title{Election-night model: specification}
\author{Young-Kyung Kim}
\date{\today}

\begin{document}
\maketitle

\section{Outcomes}

The model estimates three outcomes at the precinct level, separately by vote mode, giving
nine equations in total. Each equation is estimated separately, but the errors are
modelled jointly through a $9 \times 9$ covariance matrix $\Sigma$.

Let $m \in \{\text{mail}, \text{early}, \text{eday}\}$, and let
$v^{\dem}_{im}$, $v^{\rep}_{im}$, $v^{\oth}_{im}$ denote votes cast in mode $m$ in
precinct $i$ for the Democratic candidate, the Republican candidate, and all others.

\begin{itemize}
  \item Three equations for $\tau_{im}$, total votes cast (turnout) in the cell;
  \item Three equations for $tp_{im}$, the third-party share of the cell;
  \item Three equations for $s_{im}$, the two-party Democratic share of the cell.
\end{itemize}

\paragraph{What $\Sigma$ does.}
Each equation is estimated separately; since they share a design matrix, seemingly
unrelated regression coincides with equation-by-equation estimation, so $\Sigma$ does not
change $\hat\beta$. It acts on prediction. The aggregate is nonlinear in $(\tau, tp, s)$
jointly within a cell, so the predictive variance depends on their correlation, and
conditioning on the observed equations sharpens the unobserved ones rather than merely
rescaling them --- the conditional covariance is a Schur complement, strictly smaller
than the unconditional one whenever the cross-blocks are nonzero.

\section{The factorization}

Predicting $v^{\dem}$, $v^{\rep}$ and $v^{\oth}$ separately and summing them can produce
impossible values --- shares that do not sum to one, negative counts. Clamping such
values afterwards would spend researcher degrees of freedom we would rather not spend.
Instead we change variables. The map
\begin{equation*}
  \tau_{im} = v^{\dem}_{im} + v^{\rep}_{im} + v^{\oth}_{im}, \qquad
  tp_{im}   = \frac{v^{\oth}_{im}}{\tau_{im}}, \qquad
  s_{im}    = \frac{v^{\dem}_{im}}{v^{\dem}_{im} + v^{\rep}_{im}}
\end{equation*}
is a bijection, so the quantity we care about can be written as a function of
$(\tau, tp, s)$ and recovered exactly. Every draw is a valid election outcome by
construction, with no clamping and no renormalizing.

\section{Link functions}

The two share outcomes are bounded in $[0,1]$, so both are modelled on the logit scale:
\begin{equation*}
  y^{tp}_{im} = \logit\!\left(\frac{v^{\oth}_{im} + \tfrac{1}{2}}{\tau_{im} + 1}\right),
  \qquad
  y^{s}_{im} = \logit\!\left(\frac{v^{\dem}_{im} + \tfrac{1}{2}}
                                  {v^{\dem}_{im} + v^{\rep}_{im} + 1}\right).
\end{equation*}
The $\tfrac{1}{2}$ offsets keep the transform finite where a third-party candidate
receives no votes in a precinct, which is common.

Turnout is modelled on a log scale, because precinct-level mode-specific turnout is
right-skewed. Where the composition is disabled the outcome is simply
$y^{\tau}_{im} = \log(\tau_{im} + 0.1)$, the offset keeping cells with zero counted votes
finite. Getting this scale right mattered more than expected.

\section{Turnout: total and composition}

By default turnout is reparameterized rather than modelled as three free log-levels,
because turnout has a known ceiling and the three modes must add up. Writing $T_i$ for
the precinct total and $p_{im}$ for its split across modes,
\begin{equation}
  \theta_i = \logit\!\left(\frac{T_i}{\kappa n_i}\right), \qquad
  \eta_{im} = \log\frac{p_{im}}{p_{i,\refm}} \quad (m \neq \refm), \qquad
  p_i = \softmax(\eta_i), \qquad
  \tau_{im} = T_i\, p_{im},
  \label{eq:compose}
\end{equation}
with $\kappa = 1.25$ and $n_i$ the precinct's registrants. The fitted outcomes are
$\theta_i$ in the reference mode's slot and the $\eta_{im}$ in the other two, so the
equation count is unchanged at $1 + (M-1) + 2M = 3M$.

This buys two guarantees. First, $\sum_m p_{im} = 1$ gives $\sum_m \tau_{im} = T_i$
exactly, so a surplus in one mode can no longer inflate the precinct total. Second,
$0 < T_i < \kappa n_i$. An earlier version of the model predicted turnout at 184\% of
truth; under \eqref{eq:compose} that is arithmetically impossible rather than something
detected after the fact.

Setting $\refm = \text{early}$ is an identifiability normalization; early is chosen
because it reports first and is the best-measured mode.

\section{Partial counts}

Treating an incomplete snapshot as final, or thresholding it as ``complete enough,''
produces strange behaviour, and the problem compounds when the decision is made per
precinct-mode rather than per precinct. We therefore model partial counts as a mixture.
With $c_{im}$ the ballots already counted in a cell and $r_{im}$ those remaining,
\begin{equation}
  r_{im} = \tau_{im} - c_{im}, \qquad
  \dem_{im} = c^{\dem}_{im} + r_{im}\,(1 - tp_{im})\,s_{im}.
  \label{eq:mixture}
\end{equation}
A counted-out cell has $r_{im} \approx 0$, so its observed share is pinned automatically;
a barely-started cell is mostly predicted. No completeness classifier is required.

Partial counts also impose their own restrictions, and these are inequalities. The floor
must act on the precinct \emph{total}: flooring each cell at its own count destroys the
adding-up property the parameterization exists for. So we floor the total and hand the
remainder only to modes short of their composition-implied share, in proportion to how
short they are:
\begin{equation}
  \begin{aligned}
    T_i &\leftarrow \max\Big(T_i,\ \textstyle\sum_m c_{im}\Big),
    &\qquad
    \text{short}_{im} &= \max\big(T_i\, p_{im} - c_{im},\, 0\big), \\[4pt]
    r_{im} &= \Big(T_i - \textstyle\sum_{m'} c_{im'}\Big)
              \frac{\text{short}_{im}}{\sum_{m'} \text{short}_{im'}}.
  \end{aligned}
  \label{eq:alloc}
\end{equation}
This preserves both $\sum_m \tau_{im} = T_i$ and $\tau_{im} \geq c_{im}$.

\paragraph{Known limitation.}
$\tau_{im} = T_i\, p_{im}$ holds only when no mode has over-reported relative to its
composition share, and $T_i < \kappa n_i$ only when $\sum_m c_{im} < \kappa n_i$. Both
can fail empirically; the unconditional invariants are $\sum_m \tau_{im} = T_i$ and
$\tau_{im} \geq c_{im}$.

\section{Covariates and the county effect}

\begin{itemize}
  \item \textbf{Turnout and share:}
        $\{\log n,\ \var{turnout22},\ \var{turnout24},\ p_{\text{male}},\ p_{65+},\
        p_{\text{black}},\ p_{\text{hisp}},\ p_{\text{col}},\ \var{dem2}\}$
  \item \textbf{Third party:}
        $\{\log n,\ p_{\text{male}},\ p_{65+},\ \var{dem2},\ \var{tp\_prior\_raw}\}$ ---
        kept deliberately small, since third-party share is a weak and stable signal that
        overfits easily.
\end{itemize}

We also include a county random effect to account for dependence across precincts.
Ignoring it is not a minor misspecification: measured across nine states, county
intraclass correlations run from $0.09$ to $0.43$, which corresponds to standard errors
understated by a factor of roughly $2.4$ to $5.1$.

\section{Estimation}

Each equation takes the form
\begin{equation}
  y^{k}_{im} = x_i' \beta^{k}_{m} + u^{k}_{m,c(i)} + \epsilon^{k}_{im},
  \qquad
  u^{k}_{m,c} \sim N\!\left(0,\ \sigma_k^2 / \lambda_k\right),
  \qquad
  \lambda_k = \frac{1 - \ICC_k}{\ICC_k}.
  \label{eq:model}
\end{equation}

Covariates from the voter file and polling priors supply a fitted prior value for every
precinct and mode. We then follow the Bayesian data-augmentation approach used in the
NYC mayoral model, with the county random effect added.

Write $X$ for the covariate block and $C$ for the county indicator block. The design
stacks $N$ $g$-prior rows $(X, \mathbf{0})$ at weight $n_k / N$, the $n_{\text{obs}}$
reported rows $(X_{\text{obs}}, C_{\text{obs}})$ at weight $w_{\text{obs}}$, and $J$
ridge rows $(\mathbf{0}, I_J)$ at weight $\lambda_k$, against the response
$(\tilde{y}^{k},\ y^{k}_{\text{obs}},\ \mathbf{0})$:
\begin{equation}
  A = \begin{bmatrix}
        X & \mathbf{0} \\
        X_{\text{obs}} & C_{\text{obs}} \\
        \mathbf{0} & I_J
      \end{bmatrix},
  \qquad
  y = \begin{bmatrix} \tilde{y}^{k} \\ y^{k}_{\text{obs}} \\ \mathbf{0} \end{bmatrix},
  \qquad
  W = \operatorname{diag}\!\left(\tfrac{n_k}{N}\mathbf{1}_N,\ w_{\text{obs}},\
      \lambda_k \mathbf{1}_J\right).
  \label{eq:design}
\end{equation}

The first block is the $g$-prior, with the county columns zeroed so that it affects
$\beta$ only. The second is the reported data, with the share equations weighted by
counted ballots --- normalised to mean one, so $n_k$ keeps its interpretation in
precinct-equivalents --- and the turnout equation unweighted. The third shrinks $u$
toward zero with weight $\lambda_k$; it is these rows, not $C$ itself, that make the
county coefficients a random effect. The whole system is one weighted least squares
solve, which keeps conjugacy and closed-form draws.

\paragraph{The covariance matrix.}
$\Sigma$ is assembled as $D\,\mathcal{C}\,D$: the nine per-equation residual standard
deviations $D$ come from the current night's fits, while the correlations $\mathcal{C}$
are supplied from past elections, because the modes report in sequence and the cross-mode
blocks are never estimable from tonight's data. The measured within-mode correlation
between turnout and share varies in sign across modes, so a single pooled value will not
do.

\section{Prediction}

Take 500 posterior draws. Transform each draw into the relevant units ---
$\kappa n_i \cdot \logit^{-1}(\theta)$ for the precinct total, $\softmax$ for the mode
composition, and $\logit^{-1}$ for the two shares --- so that every draw satisfies the
adding-up and bound constraints by construction. Then combine counted and remaining votes
per cell using \eqref{eq:mixture} and \eqref{eq:alloc}, and aggregate.

\end{document}
